\section{Introduction}
\label{sec:introduction}
In robotic manipulation, deciding what to do next often requires information from past observations. A robot may need to remember which container hides a target, count earlier events, or reproduce a demonstrated action sequence. RoboMME evaluates these capabilities through 16 tasks grouped into counting, permanence, reference, and imitation \citepending{robomme}. Solving them requires both interpreting the relevant history and executing actions consistent with it.

Recent RoboMME studies use GPT-6 Astra (hereafter Astra), a vision-language reasoning model, in two distinct roles. Direct control asks Astra to generate robot motion targets from visual history, placing action generation as well as task reasoning in the model and incurring repeated inference during execution \citepending{astra-direct}. A three-tier system instead uses Astra to plan subgoals for a trained vision-language-action model (VLA), with a separate learned monitor tracking completion \citepending{astra-three-tier}. This division reduces Astra's involvement in routine motion, but concentrates its role on planning. We seek to retain the efficiency of learned control while also using Astra to supervise progress and intervene when execution goes wrong.

Reviewing every short action sequence, or action chunk, incurs repeated inference cost even when the VLA is still executing the same subgoal. Longer intervals between reviews can delay correction of a failed grasp or an incorrect motion. Appropriate supervision therefore depends on the current task stage, visual evidence, and the controller's behavior in similar situations. Experience from earlier episodes can help guide these decisions, but must be updated when new outcomes reveal its limitations.

We present a self-evolving harness in which Astra learns from execution experience how to supervise the official RoboMME GroundSG VLA (Figure~\ref{fig:teaser}). Given task history and current observations, Astra selects a subgoal and a finite execution budget. When permitting a longer interval, it can activate auxiliary monitors, such as gripper-state or directional-motion checks, to trigger review before the budget expires. The VLA continues to react to the latest observations throughout this interval. At review, Astra decides whether to continue execution, revise the subgoal, or intervene directly. Post-episode reflection revises advice about these choices, progressively informing supervision without changing VLA weights.

We evaluate the harness on all 16 RoboMME tasks with 50 episodes per task. It achieves 97.1\% success across 800 episodes, reaches at least 90\% on every task, and executes 54.9\% of control steps without additional Astra review. Execution and experience traces show how prior outcomes lead to revised review intervals and more specific task interpretations. Together, these results and traces illustrate how experience can guide Astra's supervision while the VLA carries out substantial portions of execution without additional review.
